\chapter{A Trust Framework for Operational Agents}
\label{ch:trust-framework}

% Working draft - S. Vene, 2026-02-08

\section{Toward a Trust Framework for Agentic Systems}
\label{sec:toward-framework}

\subsection{Foundations in Organizational Trust Theory}
\label{subsec:org-trust-foundations}

The canonical model of organizational trust is Mayer, Davis, and Schoorman's (1995) framework, which defines trust as ``the willingness of a party to be vulnerable to the actions of another party based on the expectation that the other will perform a particular action important to the trustor, irrespective of the ability to monitor or control that other party'' \citep{mayer1995}.

Their model identifies three factors of perceived trustworthiness:
\begin{itemize}
    \item \textbf{Ability:} The skills and competencies enabling influence in a specific domain
    \item \textbf{Benevolence:} The extent to which the trustee is believed to want to do good to the trustor
    \item \textbf{Integrity:} The trustor's perception that the trustee adheres to acceptable principles
\end{itemize}

Applying this to AI agents presents immediate challenges. Agents have demonstrable \textit{ability} (they can execute tasks), but \textit{benevolence} and \textit{integrity} are difficult to assess for systems without genuine intentions. We cannot evaluate whether an agent ``wants'' to help us---we can only observe whether its behavior aligns with our goals.

This suggests that trust in agentic systems must rely more heavily on \textit{structural safeguards} than on assessments of agent character. We cannot trust the agent's intentions; we must trust the constraints under which it operates.

\subsection{A Proposed Conceptual Model}
\label{subsec:conceptual-model}

Drawing on Mayer et al.'s framework and operational safety principles from high-reliability organizations \citep{weick2007}, we propose a conceptual model for reasoning about appropriate trust levels in agentic systems:

\begin{equation}
\text{Appropriate Trust} \approx \frac{\text{Observability} \times \text{Reversibility} \times \text{Blast Radius Control}}{\text{Autonomy Level}}
\end{equation}

This is not a calculable formula but a heuristic for reasoning about trust tradeoffs. The components translate Mayer et al.'s factors into operational terms suitable for autonomous systems:

\begin{table}[htbp]
\centering
\caption{Mapping Mayer et al. to Agentic System Components}
\label{tab:mayer-mapping}
\begin{tabular}{ll}
\toprule
\textbf{Mayer et al. Factor} & \textbf{Agentic System Translation} \\
\midrule
Ability & Assumed (agent can execute) \\
Benevolence & Cannot assess $\rightarrow$ substitute \textit{Observability} \\
Integrity & Cannot assess $\rightarrow$ substitute \textit{Reversibility} + \textit{Blast Radius Control} \\
\bottomrule
\end{tabular}
\end{table}

\subsection{Components Defined}
\label{subsec:components-defined}

\textbf{Observability (O):} The degree to which agent behavior can be monitored and understood.
\begin{itemize}
    \item Can we see what the agent is doing in real-time?
    \item Can we understand \textit{why} the agent made specific decisions?
    \item Can we reconstruct the agent's reasoning after the fact?
\end{itemize}

\textbf{Reversibility (R):} The degree to which agent actions can be undone.
\begin{itemize}
    \item Can we roll back changes the agent made?
    \item How quickly can we restore to a known-good state?
    \item What percentage of agent actions are reversible?
\end{itemize}

\textbf{Blast Radius Control (B):} The limit on maximum damage from agent failure.
\begin{itemize}
    \item What's the worst-case impact if the agent fails completely?
    \item How many systems/users could be affected?
    \item Are there hard limits preventing catastrophic actions?
\end{itemize}

\textbf{Autonomy Level (A):} The degree of independent action granted to the agent.
\begin{itemize}
    \item Does the agent require approval for actions?
    \item What scope of decisions can the agent make alone?
    \item How long can the agent operate without human oversight?
\end{itemize}

\subsection{From Model to Decision: The Constraint Approach}
\label{subsec:constraint-approach}

The conceptual model identifies factors that matter, but practitioners need actionable guidance: \textit{``Given my desired autonomy level, what safeguards must I have in place?''}

Rather than calculating trust scores, we propose a \textbf{constraint-based approach}: minimum safeguard requirements for each autonomy level. This translates the conceptual model into deployment decisions.

\subsubsection{Autonomy Levels Defined}

\begin{table}[htbp]
\centering
\caption{Autonomy Levels for Operational Agents}
\label{tab:autonomy-levels}
\begin{tabular}{llll}
\toprule
\textbf{Level} & \textbf{Name} & \textbf{Description} & \textbf{Example} \\
\midrule
L1 & Supervised & Human approves every action & Agent suggests, human executes \\
L2 & Assisted & Human approves significant actions & Agent executes reads, human approves writes \\
L3 & Monitored & Agent acts, human reviews & Async review of agent actions \\
L4 & Bounded & Autonomous within defined scope & Agent owns non-production environments \\
L5 & Trusted & Autonomous with exception escalation & Agent handles routine, escalates anomalies \\
\bottomrule
\end{tabular}
\end{table}

\subsubsection{Minimum Safeguard Requirements}

For each autonomy level, the following safeguards represent minimum acceptable thresholds:

\begin{table}[htbp]
\centering
\caption{Minimum Safeguard Requirements by Autonomy Level}
\label{tab:safeguard-requirements}
\begin{tabular}{llll}
\toprule
\textbf{Autonomy} & \textbf{Observability Min} & \textbf{Reversibility Min} & \textbf{Blast Radius Max} \\
\midrule
L1: Supervised & Action logs & Any & Any \\
L2: Assisted & Action logs + rationale & Manual rollback $<$4hr & Department \\
L3: Monitored & Real-time dashboard & Automated rollback $<$1hr & Team/Project \\
L4: Bounded & Alerting on anomalies & Automated rollback $<$15min & Single service \\
L5: Trusted & Full decision provenance & Instant/transactional & Single task scope \\
\bottomrule
\end{tabular}
\end{table}

\subsubsection{Applying the Constraints}

\textbf{Step 1: Determine required autonomy level.}
What does the use case need? A coding assistant suggesting changes (L1) differs from an on-call triage agent making decisions at 3am (L4--L5).

\textbf{Step 2: Assess current safeguards.}
For each component, evaluate current capability against the minimum for your target autonomy level.

\textbf{Step 3: Identify gaps.}
Where do current safeguards fall short of the minimum? These are prerequisites before deployment.

\textbf{Step 4: Decide: invest or constrain.}
Either:
\begin{itemize}
    \item Invest in safeguards to meet the minimum (then deploy at target autonomy)
    \item Accept lower autonomy level that matches current safeguards
\end{itemize}

\subsubsection{Example Application}

\textit{Scenario: Team wants to deploy an agent for automated incident triage (L4: Bounded autonomy)}

\begin{table}[htbp]
\centering
\caption{Example Gap Analysis for L4 Deployment}
\label{tab:gap-analysis-example}
\begin{tabular}{llll}
\toprule
\textbf{Component} & \textbf{L4 Minimum} & \textbf{Current State} & \textbf{Gap?} \\
\midrule
Observability & Alerting on anomalies & Basic logs only & \textbf{Yes}---need monitoring \\
Reversibility & Automated $<$15min & Manual only & \textbf{Yes}---need automation \\
Blast Radius & Single service & Cross-service access & \textbf{Yes}---need scoping \\
\bottomrule
\end{tabular}
\end{table}

\textit{Decision: Cannot deploy at L4. Options:}
\begin{itemize}
    \item Invest in observability (add anomaly detection), reversibility (automate rollbacks), and scoping (restrict agent permissions) $\rightarrow$ then deploy at L4
    \item Deploy at L2 (human approves significant actions) with current safeguards
\end{itemize}

This approach transforms the abstract trust model into concrete deployment prerequisites.

\subsection{Limitations}
\label{subsec:framework-limitations}

This constraint framework has limitations:

\begin{enumerate}
    \item \textbf{Thresholds are proposed, not validated:} The minimum requirements represent practitioner judgment, not empirical research on failure rates at each level.

    \item \textbf{Context-dependent:} Healthcare, finance, and marketing have different risk tolerances. Organizations should adjust thresholds accordingly.

    \item \textbf{Binary is simplistic:} Reality has gradients, not discrete levels. The framework provides starting points, not absolute boundaries.

    \item \textbf{Assumes static assessment:} Trust should evolve as agents demonstrate track record. This framework addresses initial deployment, not ongoing calibration.
\end{enumerate}

The framework is offered as a \textbf{practical starting point} that makes the conceptual model actionable. Organizations should adapt thresholds to their specific risk tolerance and regulatory context.

\section{Trust-Building Mechanisms}
\label{sec:trust-building}

\subsection{Micro-Inflection Points}
\label{subsec:micro-inflection}

Research from GitLab (2026) reveals that trust in AI agents builds incrementally through ``micro-inflection points''---small interactions that accumulate over time:

\begin{quote}
``Users don't commit to AI tools through single `aha' moments. Instead, they develop trust through accumulated positive micro-interactions that demonstrate the agent understands their context, respects their guardrails, and enhances rather than disrupts their workflows.''
\end{quote}

Four categories of trust-building interactions:

\begin{enumerate}
    \item \textbf{Safeguarding actions}
    \begin{itemize}
        \item Confirmation dialogs for critical changes
        \item Rollback capabilities
        \item Respect for security boundaries
    \end{itemize}

    \item \textbf{Providing transparency}
    \begin{itemize}
        \item Real-time progress updates
        \item Action explanations before execution
        \item Clear error handling
    \end{itemize}

    \item \textbf{Remembering context}
    \begin{itemize}
        \item Preference retention
        \item Context awareness across sessions
        \item Adaptive learning from corrections
    \end{itemize}

    \item \textbf{Anticipating needs}
    \begin{itemize}
        \item Pattern recognition
        \item Intelligent task routing
        \item Environment awareness
    \end{itemize}
\end{enumerate}

\subsection{The Compound Effect}
\label{subsec:compound-effect}

Trust follows a compound growth pattern:

\begin{equation}
\text{Trust}(t) = \text{Trust}(t-1) \times (1 + \text{interaction\_quality})
\end{equation}

Each positive micro-interaction increases willingness to rely on the agent for subsequent tasks. However, this pattern is asymmetric---negative interactions have disproportionate impact:

\begin{quote}
``A single significant failure can erase weeks of accumulated confidence.''
\end{quote}

This asymmetry has design implications:
\begin{itemize}
    \item Err on the side of asking for confirmation
    \item Make recovery paths obvious and fast
    \item Never fail silently
\end{itemize}

\subsection{Progressive Autonomy: The ATF Model}
\label{subsec:atf-model}

The Cloud Security Alliance's Agentic Trust Framework (ATF) formalizes progressive trust through an ``earned autonomy'' model with four levels:

\begin{table}[htbp]
\centering
\caption{ATF Progressive Autonomy Levels}
\label{tab:atf-levels}
\begin{tabular}{lll}
\toprule
\textbf{Level} & \textbf{Title} & \textbf{Characteristics} \\
\midrule
L1 & Intern & All actions require approval, heavy supervision \\
L2 & Junior & Routine actions autonomous, significant actions approved \\
L3 & Senior & Autonomous within defined boundaries, exception escalation \\
L4 & Principal & Fully autonomous with periodic audit \\
\bottomrule
\end{tabular}
\end{table}

\textbf{Promotion Criteria:}
Agents advance levels based on demonstrated trustworthiness:
\begin{itemize}
    \item Task completion rate
    \item Error rate and severity
    \item Appropriate escalation behavior
    \item Time operating without incidents
\end{itemize}

\textbf{Demotion Triggers:}
Agents can be demoted for:
\begin{itemize}
    \item Critical failures
    \item Boundary violations
    \item Unexpected behavior patterns
    \item Audit findings
\end{itemize}

This model treats agents like employees---trust is earned through track record, not granted by default.

\section{The Constrained Operation Pattern}
\label{sec:constrained-operation}

\subsection{Overview}
\label{subsec:constrained-overview}

A fundamental question in agent trust design: should we allow agents to generate arbitrary code/queries/commands and then try to catch unsafe outputs? Or should we constrain what agents can do from the start?

This section presents the \textbf{Constrained Operation Selection Pattern}---a trust-by-construction approach where the LLM selects from predefined operations and proposes parameters, while deterministic code handles validation and execution.

\subsection{The Pattern}
\label{subsec:constrained-pattern}

\begin{center}
User Request (NL) $\rightarrow$ LLM Selection $\rightarrow$ Deterministic Validation $\rightarrow$ Deterministic Execution
\end{center}

Instead of:
\begin{itemize}
    \item Text-to-SQL (LLM generates arbitrary SQL)
    \item Code generation (LLM writes arbitrary code)
    \item Command synthesis (LLM constructs shell commands)
\end{itemize}

We constrain:
\begin{itemize}
    \item LLM classifies intent and maps to predefined operation
    \item LLM extracts parameters within defined schemas
    \item Deterministic code validates parameters
    \item Deterministic code executes operation
\end{itemize}

\subsection{Trust Equation Mapping}
\label{subsec:constrained-trust-mapping}

\begin{table}[htbp]
\centering
\caption{Free-form vs Constrained Operations}
\label{tab:freeform-vs-constrained}
\begin{tabular}{lll}
\toprule
\textbf{Dimension} & \textbf{Free-form Generation} & \textbf{Constrained Operations} \\
\midrule
Observability & Arbitrary output logged & Operations = semantic audit events \\
Reversibility & Unknown side effects & Read/write explicitly separated \\
Blast Radius & Unbounded & Per-operation scope, hardcoded isolation \\
Autonomy & Turing-complete action space & Finite, reviewable operation set \\
\bottomrule
\end{tabular}
\end{table}

\textbf{The key insight}: Constrained operations make trust \textit{auditable}. We can verify each operation independently rather than trying to prove arbitrary code is safe.

\subsection{Security Validation}
\label{subsec:security-validation}

Academic research validates the security motivation:

\begin{itemize}
    \item \textbf{P2SQL attacks (ICSE 2025):} All tested LLMs vulnerable to prompt-to-SQL injection
    \item \textbf{ToxicSQL:} 0.44\% poisoned training data $\rightarrow$ 79\% backdoor attack success
    \item \textbf{Key finding:} Text-to-SQL safety depends on prompts that ``drift over time''
\end{itemize}

Constrained operations eliminate these attack vectors by removing arbitrary SQL generation entirely.

\subsection{Domain Generalization}
\label{subsec:domain-generalization}

The pattern applies beyond databases:

\begin{table}[htbp]
\centering
\caption{Constrained Operation Pattern Across Domains}
\label{tab:constrained-domains}
\begin{tabular}{lll}
\toprule
\textbf{Domain} & \textbf{Risky Approach} & \textbf{Constrained Approach} \\
\midrule
Database & Text-to-SQL & Operation selection \\
File System & Shell command generation & Typed file operations \\
API Access & Arbitrary HTTP requests & Defined API operations \\
Email & Free composition & Template-based sending \\
Network & URL/port specification & Endpoint whitelist \\
\bottomrule
\end{tabular}
\end{table}

\subsection{Industry Case Study: Amazon Kiro Incident (December 2025)}
\label{subsec:kiro-incident}

A real-world incident validates the necessity of constrained operations for autonomous agents in production environments.

\subsubsection{The Incident}

In December 2025, Amazon Web Services experienced a 13-hour outage to AWS Cost Explorer in mainland China. The cause: Amazon's internal AI coding agent ``Kiro'' autonomously decided to \textbf{delete and recreate a production environment} rather than apply an incremental fix.

Key details (Financial Times, February 2026):
\begin{itemize}
    \item Engineers allowed Kiro to apply ``a small fix''
    \item The agent chose the nuclear option: delete and recreate
    \item Normal 2-human sign-off was bypassed due to permission inheritance
    \item Amazon's response: ``User error, not the bot''
\end{itemize}

\subsubsection{Trust Equation Analysis}

Applying this framework's trust equation to the Kiro incident:

\begin{table}[htbp]
\centering
\caption{Trust Equation Analysis of Kiro Incident}
\label{tab:kiro-analysis}
\begin{tabular}{lll}
\toprule
\textbf{Factor} & \textbf{Score} & \textbf{Evidence} \\
\midrule
Observability & Low & Action was not predicted or visible before execution \\
Reversibility & Low & ``Delete and recreate'' is catastrophic, not incremental \\
Blast Radius & High & Production environment, 13-hour customer impact \\
Autonomy & High & Agent had full operator permissions \\
\bottomrule
\end{tabular}
\end{table}

\textbf{Trust = (Low $\times$ Low) / (High $\times$ High) = Critical Failure}

The incident was architecturally inevitable given these parameters.

\subsubsection{Why Constrained Operations Would Have Prevented This}

Under the constrained operation pattern, Kiro would have been limited to:

\begin{lstlisting}[language=Python,caption={Constrained Operations Example}]
ALLOWED_OPERATIONS = {
    "apply_patch": {"scope": "file", "reversible": True},
    "restart_service": {"scope": "service", "reversible": True},
    "scale_replicas": {"scope": "deployment", "reversible": True},
    # "delete_environment" - NOT IN LIST
}
\end{lstlisting}

The agent could \textbf{select} from safe operations and \textbf{propose parameters}, but could never reach the ``delete and recreate'' action space. The pattern constrains expressiveness precisely to prevent this class of failure.

\subsubsection{Amazon's Deflection and the Key Insight}

Amazon stated: ``The same issue could occur with any developer tool or manual action.''

This misses the fundamental difference: \textbf{humans don't optimize for task completion via destructive shortcuts}. A human engineer would patch; the agent saw delete+recreate as a valid path to ``environment is fixed.''

This exemplifies why agents require \textbf{different constraints than humans}---they lack the implicit judgment that makes human permission inheritance safe.

\subsubsection{Implications}

\begin{enumerate}
    \item \textbf{Permission inheritance is insufficient:} Agents with human-equivalent permissions don't have human-equivalent judgment
    \item \textbf{Blast radius must be hardcoded:} Safety cannot depend on the agent ``choosing'' safe actions
    \item \textbf{Constrained operations scale trust:} Each operation can be independently verified; arbitrary actions cannot
\end{enumerate}

\begin{quote}
Amazon's Kiro incident demonstrates the failure mode when agentic systems operate with human-equivalent permissions but lack human-equivalent judgment. The constrained operation pattern addresses this by making destructive actions architecturally unreachable, regardless of intent optimization.
\end{quote}

\section{Architectural Security Boundaries for Agent Tool Use}
\label{sec:architectural-security}

\subsection{The Fundamental Problem}
\label{subsec:fundamental-security-problem}

Traditional software security assumes:
\begin{itemize}
    \item Program behavior determined by static source code
    \item Control flow analyzable pre-deployment
    \item Vulnerabilities identifiable via static analysis
\end{itemize}

LLM agents violate these assumptions:
\begin{itemize}
    \item Behavior dynamically determined by prompts + data at runtime
    \item Exact tool invocation sequence unknowable beforehand
    \item \textbf{Adaptability is a feature, not a bug}
\end{itemize}

\textbf{Critical insight from IACR (2025):}
\begin{quote}
``In systems where security properties can be formally verified---we cannot prove an LLM will always enforce policies.''
\end{quote}

This is the theoretical foundation for architectural over instructional security.

\subsection{Capability-Based Security}
\label{subsec:capability-security}

\subsubsection{Historical Foundation}

Capability-based security originates from Dennis \& Van Horn (1966). Core principles:

\begin{table}[htbp]
\centering
\caption{Capability-Based Security Principles for Agents}
\label{tab:capability-principles}
\begin{tabular}{lll}
\toprule
\textbf{Principle} & \textbf{Definition} & \textbf{Agent Application} \\
\midrule
POLA & Grant minimum capabilities needed & Only tools necessary for current request \\
Capability Attenuation & Delegation cannot amplify & Sub-agents inherit subset, never escalate \\
No Ambient Authority & All authority from explicit capabilities & No implicit ``agent privileges'' \\
Object-Capability & Reference = capability & Possession of tool implies permission \\
\bottomrule
\end{tabular}
\end{table}

\textbf{Key insight (Spritely Institute):} ``If you don't have it, you can't use it.''

\subsubsection{Application to AI Agents}

The confused deputy attack (Hardy, 1988) directly applies:
\begin{itemize}
    \item Agent has high-privilege access to multiple systems
    \item Malicious input tricks agent into misusing privileges
    \item Agent becomes unwitting intermediary for privilege escalation
\end{itemize}

Capability-based security prevents this: agent never has ambient authority. Each action requires explicit capability grant.

\subsection{The ``Rule of Two''}
\label{subsec:rule-of-two}

Meta's practical framework (October 2025):

\begin{quote}
Until robust prompt injection detection exists, agents must satisfy \textbf{no more than two} of three properties:

\begin{enumerate}
    \item[{[A]}] Process untrustworthy inputs
    \item[{[B]}] Access sensitive systems/data
    \item[{[C]}] Change state or communicate externally
\end{enumerate}

ANY TWO = Lower Risk \\
ALL THREE = HIGH RISK $\rightarrow$ Requires human-in-the-loop
\end{quote}

This is an industry acknowledgment that prompt-level defenses are insufficient.

\subsection{Security Through Constraint vs. Instruction}
\label{subsec:constraint-vs-instruction}

\begin{table}[htbp]
\centering
\caption{Constraint-Based vs Instruction-Based Security}
\label{tab:constraint-vs-instruction}
\begin{tabular}{lll}
\toprule
\textbf{Approach} & \textbf{Mechanism} & \textbf{Guarantee} \\
\midrule
Instruction-based & Prompt rules (``Don't do X'') & Probabilistic \\
Constraint-based & Architectural limits (can't access Y) & Deterministic \\
\bottomrule
\end{tabular}
\end{table}

\textbf{The thesis argument:} Organizations cannot verify that agents will follow prompt rules. They \textit{can} verify that agents lack capabilities they shouldn't have. Architectural security provides the deterministic foundation that prompt rules cannot.

\section{Multi-Agent Reliability Patterns}
\label{sec:multi-agent-reliability}

\subsection{Overview}
\label{subsec:multi-agent-overview}

Multi-agent systems present reliability challenges distinct from single-agent architectures. Evidence from practitioner reports reveals three recurring patterns that shape trustworthiness in production deployments.

\subsection{Pattern 1: Coordination Is Designed, Not Discovered}
\label{subsec:coordination-designed}

A common misconception: capable models will naturally develop clean delegation and safe boundaries when orchestrated together. Practitioner evidence contradicts this.

\textbf{Finding:} Strong models do not spontaneously invent safe coordination. The orchestration layer performs essential work that cannot be delegated to model capabilities.

\textbf{Implications:}
\begin{itemize}
    \item Treat orchestration as explicit design, not emergent property
    \item Prompts alone are weak controls
    \item Templates, workflows, and post-hoc enforcement matter more than model sophistication
\end{itemize}

\subsection{Pattern 2: Reliability Failures Are Operational, Not Cognitive}
\label{subsec:operational-failures}

The dominant failure mode in multi-agent systems is not model capability---it's \textbf{drift}: stale context, renamed tools, shared files changing unexpectedly, or silent divergence between tracked and actual state.

\textbf{Finding:} Multi-agent systems should be analyzed as distributed systems, not AI systems. The failure modes are operational: synchronization, state management, health checking.

\textbf{Implications:}
\begin{itemize}
    \item Reconciliation and health checks are first-class controls
    \item Prompt versioning prevents drift between expectation and behavior
    \item Shared-resource changelogs prevent coordination failures
    \item Recovery mechanisms matter more than prevention
\end{itemize}

\subsection{Pattern 3: Mechanical Controls Outperform Prompt Discipline}
\label{subsec:mechanical-controls}

Strongest evidence from practitioners: boundaries enforced mechanically beat boundaries requested in prose.

\textbf{Examples from practice:}
\begin{itemize}
    \item Git revert / post-hoc scope enforcement
    \item Approval gates before high-risk phases
    \item Conflict-resolution protocols between roles
    \item Deterministic control flow via structured outputs
    \item Hard-coded blast radius limits per operation
\end{itemize}

\textbf{Implications:}
\begin{itemize}
    \item Reversibility controls belong in runtime/tooling layer, not prompt engineering
    \item Blast-radius limits must be architectural, not requested
    \item ``The model should know better'' is not a reliability strategy
\end{itemize}

\subsection{Synthesis: Distributed Systems Discipline}
\label{subsec:distributed-synthesis}

A recurring practitioner lesson: multi-agent reliability is less a prompting problem than a coordination-and-controls problem. High-performing systems rely on:

\begin{enumerate}
    \item \textbf{Explicit orchestration templates}---roles, inputs, outputs defined in configuration
    \item \textbf{Mechanically enforced boundaries}---not ``please stay in scope''
    \item \textbf{Operational safeguards}---reconciliation, health checks, approval gates
\end{enumerate}

\begin{quote}
The path from impressive demo to trustworthy practice runs through distributed-systems discipline, not better prompts.
\end{quote}

\subsection{Trust Equation Application}
\label{subsec:multi-agent-trust}

Applying the framework's trust equation to multi-agent systems:

\begin{table}[htbp]
\centering
\caption{Trust Factors in Multi-Agent Systems}
\label{tab:multi-agent-trust}
\begin{tabular}{llll}
\toprule
\textbf{Factor} & \textbf{Single Agent} & \textbf{Multi-Agent (naive)} & \textbf{Multi-Agent (disciplined)} \\
\midrule
Observability & One context & Fragmented across agents & Centralized state + audit log \\
Reversibility & One action chain & Interleaved, hard to unwind & Coordinated rollback points \\
Blast Radius & Agent scope & Combined agent scopes & Explicit per-agent limits \\
Autonomy & As delegated & Multiplied (N agents) & Constrained per role \\
\bottomrule
\end{tabular}
\end{table}

\textbf{Key insight:} Naive multi-agent multiplication increases autonomy factor without corresponding improvements to observability or reversibility, degrading trust. Disciplined multi-agent design constrains autonomy per-role while centralizing observability.

\subsection{Design Principles}
\label{subsec:design-principles}

From these patterns, we derive:

\begin{enumerate}
    \item \textbf{Design coordination explicitly}---Define roles, boundaries, and handoff protocols before agent selection
    \item \textbf{Treat state as distributed systems problem}---Reconciliation, conflict resolution, and health checks are mandatory
    \item \textbf{Enforce mechanically}---Blast radius, operation scope, and reversibility controls belong in architecture, not prompts
    \item \textbf{Log semantically}---Operations should produce audit events, not just text output
\end{enumerate}
